%% ma: L12-B19-0030
%% lop: 12
%% chuong: 6
%% bai: 19
%% dang: 12.19.4
%% loai: TLN
%% muc: VD
%% dap_an: "13"
%% nguon: "(NBV)-10 ĐỀ KIỂM TRA HK2-TOÁN 12-DUNG CHUNG 3 SÁCH.pdf (H02, Đề số 2), Phần 3 Câu 1"
%% kien_thuc: "Bài 19 (công thức xác suất toàn phần, lấy không hoàn lại)"
%% trang_thai: chinh-thuc
%% ghi_chu: "Đáp án gốc 13 khớp. Bỏ ảnh phòng thi (ảnh trang trí)."
\begin{ex}
Một cuộc thi khoa học có $36$ bộ câu hỏi, trong đó có $20$ bộ câu hỏi về chủ đề tự nhiên và $16$ bộ câu hỏi về chủ đề xã hội. Bạn An lấy ngẫu nhiên $1$ bộ câu hỏi (lấy không hoàn lại), sau đó bạn Bình lấy ngẫu nhiên $1$ bộ câu hỏi. Xác suất bạn Bình lấy được bộ câu hỏi về chủ đề xã hội bằng $\dfrac ab$ với $\dfrac ab$ là phân số tối giản. Giá trị $a+b$ bằng bao nhiêu?
\shortans{13}
\loigiai{Gọi $A$: ``An lấy được bộ câu hỏi tự nhiên'', $B$: ``Bình lấy được bộ câu hỏi xã hội''. Khi đó $P(A)=\dfrac{20}{36}=\dfrac59$, $P(\overline A)=\dfrac49$, $P(B\mid A)=\dfrac{16}{35}$, $P(B\mid\overline A)=\dfrac{15}{35}$.

$P(B)=\dfrac59\cdot\dfrac{16}{35}+\dfrac49\cdot\dfrac{15}{35}=\dfrac{140}{315}=\dfrac49$. Vậy $a=4$, $b=9$ và $a+b=13$.}
\end{ex}
